Sean \(T\) la separación entre los dos primeros máximos y \(X_1,X_2\) sus alturas verdaderas. Con \(d=\ln(X_1/X_2)\),
\[
\alpha=\frac dT,\qquad \omega_d=\frac{2\pi}{T},\qquad
\omega_0=\frac{\sqrt{4\pi^2+d^2}}{T}>\omega_d.
\]
Al redondear hacia fuera, las proyecciones de los parámetros compatibles quedan contenidas en estas cotas, en \(\mathrm s^{-1}\):
\[
\alpha\in[0.323337,0.370280],\quad
\omega_d\in[3.110487,3.173326],\quad
\omega_0\in[3.127248,3.194856].
\]
Son intervalos de cada parámetro por separado, no tres elecciones independientes. El tercer máximo debería ocurrir en \([3.96,4.04]\) s, incompatible con \([4.08,4.12]\) s. Su altura prevista pertenece a \([0.235392,0.265409]\) cm, compatible con \([0.230,0.270]\) cm. Ningún modelo de coeficientes constantes satisface todas las cotas.
